% Generated by tools/paper/render.py from main_text.py; edit the source there, not this file.
\section{Introduction}

Most analyses that compare DNA sequences begin by finding seeds, exact matches that indicate where a query and a reference may align. Read mappers chain seeds before alignment \citep{Li2013,Li2018}, whole-genome aligners anchor on matches unique to both genomes \citep{Delcher1999,Marcais2018}, and pangenome tools summarize shared sequence through matches common to many genomes \citep{Shivakumar2025}. Two families of seeds are in wide use. \emph{Sampled k-mers} retain a content-dependent subset of k-mers, so that copies of a region retain the same ones: minimizers, the k-mer of smallest hash in each window of w consecutive k-mers \citep{Schleimer2003,Roberts2004}, and syncmers, k-mers whose smallest s-mer lies at a fixed offset \citep{Edgar2021}. \emph{Maximal exact matches} (MEMs) extend until the sequences differ; maximal unique matches (MUMs) anchor whole-genome alignment \citep{Delcher1999}, and super-maximal exact matches (SMEMs), the MEMs of a read not contained in a longer one, seed BWA-MEM \citep{Li2012,Li2013}.

Seeding is a substantial cost. On short reads, single-threaded BWA-MEM spends 40--60\% of its time finding SMEMs and locating their occurrences \citep{Vasimuddin2019}, and dedicated hardware has been proposed to accelerate this step \citep{Subramaniyan2021}. The inputs keep growing, from complete telomere-to-telomere assemblies \citep{Nurk2022} to pangenomes of dozens of individuals \citep{Liao2023}, while the throughput of a processor core grows mainly through wider vector units and greater memory-level parallelism.

Each kind of seed is computed by a different tool on a different data structure. Minimizers are computed within read mappers such as minimap2 \citep{Li2018}, syncmers and strobemers within strobealign \citep{Sahlin2021,Sahlin2022}, and simd-minimizers computes random minimizers with SIMD instructions \citep{GrootKoerkamp2025simd}. MEMs and MUMs between genomes are found with suffix trees and suffix arrays in MUMmer \citep{Kurtz2004,Marcais2018}, with sparse suffix arrays \citep{Khan2009,Vyverman2013} and with hash tables of sampled k-mers in E-MEM and copMEM \citep{Khiste2015,Grabowski2019,Grabowski2023}. SMEMs are found with FM-indexes \citep{Ferragina2000} in BWA \citep{Li2013}, bwa-mem2 \citep{Vasimuddin2019}, ropebwt3 \citep{Li2024} and minibwa \citep{Li2026}. Newer seed types are supported by few tools. Mumemto relaxes multi-MUMs to matches that occur at most a given number of times, using the suffix array, LCP array and document array of the whole collection \citep{Shivakumar2025}. f-SMEMs, the maximal query intervals that occur at least f times, were introduced as k-MEMs \citep{Tatarnikov2022} and f-MEMs \citep{Gagie2026}, and ropebwt3 and minibwa compute them for read sets \citep{Li2024,Li2026}. Most of these tools parallelize across query sequences and process each sequence on one core, as a program built on a library would. We therefore compare the tools primarily on a single core, and also report ten threads, which matter most when a single long query, such as a chromosome, cannot be divided across sequences.

Here we present sporos, a library that computes all of these seeds exactly: minimizers and syncmers, a seed index for lookups, MEMs, MUMs and frequency-capped MEMs between genomes, and SMEMs and f-SMEMs of reads. It reproduces the definitions and output formats of the established tools, and on a single core it was faster than each of them. Our contributions are:

\begin{enumerate}
  \item \textbf{Frequency-capped MEMs without a collection index} (Section 3.4). The occurrences of a MEM's string in the reference correspond one to one to the MEMs whose query interval contains it (Lemma 1), so a sweep that retains the N largest interval ends counts occurrences up to a cap N.
  \item \textbf{Duplicate-free MEM enumeration} (Section 3.3). With the sparse sampling of copMEM \citep{Grabowski2019}, each MEM is reported by exactly one anchor, and anchor probes stream through a prefetching pipeline.
  \item \textbf{SMEM search with direct extension of unique matches} (Section 3.5). A match that occurs once is located and extended by comparing 32 bases at a time against the reference, and 16 reads are searched in interleaved fashion so that their memory accesses overlap.
  \item \textbf{Vectorized k-mer sampling} (Section 3.2). Minimizers and syncmers are computed on 8 or 16 SIMD lanes with exact leftmost tie-breaking and returned in sequence order from any number of threads; a seed index answers batched lookups through one prefetching pipeline.
  \item \textbf{A tested library} (Fig.~\ref{fig:overview}). One implementation per seed type serves the Mojo, C, Python and command-line interfaces; outputs equal those of MUMmer4 and copMEM2, and those of BWA except at ambiguous reference bases, and every kernel is tested against a brute-force oracle.
\end{enumerate}

On a single core of an Intel Skylake-X processor, sporos was faster than every other tool for every seed type we measured (Fig.~\ref{fig:speedups}), from 1.8 times faster than simd-minimizers for minimizers to 176 times faster than Mumemto for frequency-capped MEMs.
